\chapter{GERAÇÃO DE CÓDIGO INTERMEDIÁRIO (TAC)}

O Código Intermediário de Três Endereços (\textit{Three-Address Code} -- TAC) é a representação intermediária linear de baixo nível gerada a partir da Árvore de Sintaxe Abstrata (AST), aproximando a semântica da linguagem MiniPar das linguagens de montagem (Assembly) e facilitando fases posteriores de análise e emissão de código de máquina.

No compilador MiniPar 2026.2, o gerador de TAC está implementado em \texttt{src/tac.py} e encapsulado no componente \texttt{C\_TAC} (\texttt{src/componente\_tac.py}), satisfazendo aos contratos da interface \texttt{ITAC}.

\section{Estrutura do TAC e Linearização de Controle}

Cada instrução de três endereços possui, no máximo, um operador no lado direito e até três operandos (dois de entrada e um de destino):
\begin{itemize}
    \item Registradores Temporários: expressões complexas são decompostas recursivamente por meio da alocação de variáveis temporárias com prefixo \texttt{t} (\texttt{t0}, \texttt{t1}, \texttt{t2}, \dots).
    \item Rótulos e Saltos de Fluxo: estruturas de controle e laços (\texttt{if}, \texttt{while}, \texttt{for}) são linearizadas em desvios condicionais (\texttt{ifFalse cond goto LABEL}) e incondicionais (\texttt{goto LABEL}).
    \item Mapeamento de Paralelismo (\texttt{fork} e \texttt{join}):
    a emissão de TAC suporta as diretivas de paralelismo:
    \begin{lstlisting}[language=bash]
t0 = FORK FORK_START_0, core 2
goto ENDFORK_1
FORK_START_0:
# Corpo da subtarefa alocada ao nucleo 2
...
ENDFORK_1:
JOIN t0
    \end{lstlisting}
\end{itemize}

\section{Validação Formal do TAC no Teste 7 (Quicksort Iterativo)}

A especificação oficial do projeto estabelece a validação formal do TAC gerado para o algoritmo de Quicksort Iterativo Não-Recursivo com pilha explícita.

No algoritmo, o particionamento in-place e a pilha de limites de subvetores utilizam explicitamente um vetor manual \texttt{S} e um ponteiro de topo \texttt{sp}. O Listing \ref{lst:tac_quicksort} apresenta o trecho do TAC gerado pelo componente \texttt{C\_TAC} para o Teste 7:

\begin{lstlisting}[language=bash, caption={Código Intermediário (TAC) gerado para o Quicksort Iterativo (Teste 7)}, label={lst:tac_quicksort}]
t0 = 0
t1 = 4
S[0] = t0
S[1] = t1
sp = 2

LOOP:
if sp == 0 goto END
sp = sp - 1
t2 = S[sp]
sp = sp - 1
t3 = S[sp]
if t3 >= t2 goto LOOP

t4 = A[t2]
t5 = t3 - 1
t6 = t3

PARTITION_LOOP:
if t6 >= t2 goto PARTITION_END
t7 = A[t6]
if t7 <= t4 goto SWAP
t6 = t6 + 1
goto PARTITION_LOOP

SWAP:
t5 = t5 + 1
t8 = A[t5]
A[t5] = A[t6]
A[t6] = t8
t6 = t6 + 1
goto PARTITION_LOOP

PARTITION_END:
t5 = t5 + 1
t9 = A[t5]
A[t5] = A[t2]
A[t2] = t9

t10 = t5 - 1
S[sp] = t3
sp = sp + 1
S[sp] = t10
sp = sp + 1

t11 = t5 + 1
S[sp] = t11
sp = sp + 1
S[sp] = t2
sp = sp + 1
goto LOOP

END:
# Array A[0..4] devidamente ordenado
\end{lstlisting}
